\documentclass[12pt]{article}
\usepackage[utf8]{inputenc}
\usepackage[russian]{babel}
\usepackage{amsmath}
\usepackage{amssymb}
\usepackage{geometry}
\geometry{a4paper, margin=2cm}

\begin{document}

\section*{Задача 1.3}

\textbf{Условие:} Написать программу прибавления $3$ к двоичному
числу, записанному наоборот. Провести анализ в среднем и в худшем
случаях.

\section*{Программа}

\begin{verbatim}
alphabet = [#, 0, 1]

tapes = [ right ]

START:
    # -> #R Check_1
Check_1:
    # -> #R STOP
    0 -> 1R Inc
    1 -> 0R Check_2
Check_2:
    0 -> 0R Inc
    1 -> 1R Inc
    # -> 0R Inc
Inc:
    0 -> 1H STOP
    1 -> 0R Inc
    # -> 1H STOP
\end{verbatim}

\subsection*{Сложность}

На пустом слове:
\[
T(\varnothing)=2
\]

Дальше рассматриваем непустое слово $w$.

\paragraph{$w_1=0$.}

Сначала из $START$ в $Check_1$, затем в $Inc$. Если $w=*1^i0*...*$, то $Inc$ обрабатывает $i$ единиц и затем заменяет ноль на единицу за $i+1$ шагов. Итого:
\[
T(w)=1+1+(i+1)=i+3
\]

\paragraph{$w_1=1$.}

Сначала из $START$ в $Check_1$, затем в $Check_2$ и потом в $Inc$. 
Пусть $w = **1^j0*...*$. В $Inc$ обрабатываются эти единицы и затем
заменяет первый ноль или символ $\#$ на единицу за $j+1$ шагов. Итого:
\[
T(w)=1+1+1+(j+1)=j+4
\]

Слова длины 1:
\[
T(\#0\#)=3, T(\#1\#)=4
\]

\subsection*{Сложность в худшем}

Если $w_1=0$, то $i\leq n-1$, поэтому
\[
T(w)=i+3\leq n+2
\]

Если $w_1=1$, то $j\leq n-2$, поэтому
\[
T(w)=j+4\leq n+2
\]

Для $n\geq 2$ равенство есть на словах $w=01^{n-1}$ или $w=1^n$ - это худшие случаи.

В итоге, сложность в худшем случае:
\[
T_{\max}(n)=
\begin{cases}
2, & n=0\\
4, & n=1\\
n+2, & n\geq 2
\end{cases}
\]

\subsection*{Сложность в среднем}

Для $n=1$:
\[
T_{avg}(1)
=\frac{T(\#0\#)+T(\#1\#)}{2}=\frac{3+4}{2}
=\frac{7}{2}
\]

При $n=2$ имеются слова $01$ и $11$. Для них $T=4$, поэтому
\[
T_{avg}(2)=\frac{4+4}{2}=4
\]

Пусть дальше $n\geq 3$. Если $w_1=0$, то $T=i+3$, где $i$ -
количество единиц подряд после первого символа. Значения
$i=0,\ldots,n-3$ встречаются с вероятностями
$1/2^{i+1}$, а значение $i=n-1$ - с вероятностью
$1/2^{n-2}$. Поэтому
\[
T_{avg,0}(n)
=\sum_{i=0}^{n-3}\frac{i+3}{2^{i+1}}
  +\frac{n+2}{2^{n-2}}
=4-\frac{n+2}{2^{n-2}}
  +\frac{n+2}{2^{n-2}}=4
\]

Если $w_1=1$, то $T=j+4$, где $j$ -
количество единиц подряд после второго символа.
Значения $j=0,\ldots,n-4$ встречаются с вероятностями
$1/2^{j+1}$, а значение $j=n-2$ - с вероятностью
$1/2^{n-3}$. Поэтому
\[
T_{avg,1}(n)
=\sum_{j=0}^{n-4}\frac{j+4}{2^{j+1}}
  +\frac{n+2}{2^{n-3}}
=5-\frac{n+2}{2^{n-3}}
  +\frac{n+2}{2^{n-3}}
=5
\]

Что 0, что 1 встречаются с вероятностью $1/2$, поэтому
\[
T_{avg}(n)
=\frac{T_{avg,0}(n)
       +T_{avg,1}(n)}{2}
=\frac{4+5}{2}
=\frac{9}{2}, n\geq 3
\]

Значит, сложность в среднем:
\[
T_{avg}(n)=
\begin{cases}
2, & n=0\\
\dfrac{7}{2}, & n=1\\
4, & n=2\\
\dfrac{9}{2}, & n\geq 3
\end{cases}
\]

\end{document}